% Source: https://artofproblemsolving.com/community/c3h2262529_modular_inverses_proof_that_is_very_hard_for_me
% Stable identity: aops-topic-2262529-q01
\textbf{(a)} We have $ab\equiv1$ and $ac\equiv1\pmod m$. Then
\[b\equiv b\cdot1\equiv b(ac)=(ba)c\equiv1\cdot c=c\pmod m.\]
So any two inverses of $a$ are congruent modulo $m$.

\textbf{(b)} Indeed $1\cdot1\equiv1$ and $1\cdot5\equiv1\pmod4$, but $5\equiv1\pmod4$: the two ``inverses'' are the same residue class. The statement ``at most one inverse'' means one inverse modulo $m$, that is, one residue class; every integer congruent to an inverse is again an inverse. Among $0,1,\dots,m-1$ there is at most one inverse of $a$.
